\documentclass[12pt,a4paper]{article}

\usepackage[T1]{fontenc}
\usepackage[utf8]{inputenc}
\usepackage{amsmath,amssymb,amsthm}
\usepackage{graphicx}
\usepackage{booktabs}
\usepackage{xcolor}
\usepackage[hidelinks]{hyperref}
\usepackage[numbers]{natbib}

\newcommand{\todo}[1]{\textcolor{red}{[#1]}}

\title{Benchmarking graph AE methods for building embedding spaces for Pseudo-Time and Optimal Transport}
\author{Vladimir Opait}
\date{\today}

\begin{document}

\maketitle

\section{Methods}

\subsection{Dataset}
We use the in vitro hematopoietic differentiation time course of LARRY \cite{weinreb2020lineage}.
Mouse progenitors were barcoded with the LARRY lentiviral library and profiled by inDrops at
days 2, 4 and 6 of culture, starting from two populations: Lin$^-$Kit$^+$Sca1$^+$ (LSK) and
Lin$^-$Kit$^+$Sca1$^-$ (LK). The data contain 130{,}887 cells and 25{,}289 genes in 35 sequencing
libraries. We use the total-count normalized matrix provided by the authors. Cells carry one of
eleven type annotations; 72{,}193 are undifferentiated, and the largest mature types are
neutrophils (22{,}237) and monocytes (19{,}182).

Clonal barcodes are available for 49{,}302 cells in 5{,}864 clones (median size 4, maximum 177).
Of these clones, 3{,}894 have cells at more than one day, which gives 208{,}788 pairs of
sister cells sampled at different days. Cells of a clone are not synchronous in differentiation,
so a clone does not define a reference time for each cell, but it fixes which cell was sampled
earlier.

\paragraph{Batch structure.}
Each sequencing run corresponds to a specific day and a starting population, LK or LSK. Any
batch correction method would inevitably remove the signal of time and population as well, so
we chose to measure the size of the batch effect and keep it in mind when interpreting the
results.

The combination of starting population, day of collection, well and cell type, a stratum, gives
cells that are very close in biology, so differences between libraries within a stratum are
mostly technical. We used strata with at least 300 cells from at least two libraries (17 for LSK,
24 for LK). We computed 15 nearest neighbours in our PCA space, which is used throughout the
experiment and is described in more detail below. For each stratum we compared the share of
neighbours from the same library with the share expected under random mixing. A ratio of 1 means
no library structure. The median ratio over strata is 1.03 for LSK (range 0.91--1.09, almost
perfectly mixed) and 1.45 for LK (range 1.07--1.79, a batch effect is present).

On top of these strata we computed an ANOVA of the principal components with library as the
factor (variance explained by library, $\eta^2$, averaged over components with weights equal to
their variance). Library explains 0.2\% of the variance for LSK and 4.6\% for LK.

The LARRY data also let us measure the distance in this space between sister cells of the
same day, well and cell type that were sequenced in different libraries. Their biology should be
as close as possible, so the extra distance compared with sisters from the same library is
attributed to the batch. As a reference, each sister pair was matched with a random pair. We kept
one cell of the pair and replaced its sister with a random cell from another clone in the same
stratum. If the sisters came from one library, the random cell was taken from that library too.
If they came from different libraries, the random cell was taken from a library other than that
of the kept cell. In PCA space, sisters from different libraries are further apart than sisters
from the same library by 0 to 0.5 in median distance, depending on population and day (the
largest difference is for LK at day 4). The median distance between sisters is 5--10, and
between a cell and the random cell 9--14.

We conclude that there is no critical batch effect for the LSK population, and for LK it is
intermediate and we must accept it as a limitation when we will interpret the result

\subsection{Evaluation}
The pseudotime of a cell is compared with the time order within its clone. For a pair of sister
cells from the same clone, sampled at different days and assigned to the same branch, the pair
is correct if the cell from the later day has the larger pseudotime (ties count as 0.5). The
accuracy \emph{acc} is the share of correct pairs. To measure how large the errors are, the
pseudotime is converted to percentile ranks within the branch, and \emph{violation} is the mean
of $\max(0, r_{\mathrm{early}} - r_{\mathrm{late}})$ over pairs. A random pseudotime gives
acc $=0.5$ and violation $=1/6$. Confidence intervals come from a bootstrap over clones, not over
cells or pairs.

Branches are defined by the trajectory files of the original study. The neutrophil branch
contains the 61{,}310 cells of their neutrophil trajectory (63{,}079 sister pairs); the monocyte
branch contains the 36{,}148 remaining cells of their neutrophil--monocyte trajectory
(34{,}089 pairs). Metrics are computed per branch and averaged with equal weight; a branch needs
at least 200 pairs to enter the mean. As a check of the metric, the original pseudotime of the
neutrophil trajectory gives acc $=0.940$ (95\% CI 0.924--0.956) and violation $=0.004$, and a
random pseudotime gives acc $=0.501$ and violation $=0.167$.

\todo{evaluation for Waddington-OT: decided (W1 for held-out day 4 after GPA alignment, as in Zhou 2026; see claude.md, Metrics, item 4), to be written up}

\subsection{Trajectory methods}
For each method we describe what it uses from the latent space. The assumptions for the latent
spaces in Section~\ref{sec:latent} refer to these properties.

\paragraph{Euclidean distance.}
The pseudotime of a cell is its Euclidean distance to the root in the latent space. No graph is
used. This baseline shows how much of the order comes from the latent space alone.
\todo{root as for DPT}

\emph{What it uses.} \todo{which properties of the latent space this baseline relies on}

\paragraph{DPT.}
Diffusion pseudotime (DPT) \cite{haghverdi2016diffusion} was computed with Scanpy on all cells.
A $k$-nearest-neighbour graph was built on the latent space under evaluation, and DPT
distance was computed from 10 diffusion components.

\emph{What it uses.} \todo{which properties of the latent space DPT relies on}

\paragraph{Slingshot.}
\todo{description}

\emph{What it uses.} \todo{which properties of the latent space Slingshot relies on}

\paragraph{Palantir.}
Palantir \cite{setty2019palantir} computes pseudotime as the shortest-path distance from the root
on a $k$-nearest-neighbour graph built on diffusion components, and also gives fate probabilities.
\todo{parameters; how the latent space is passed in}

\emph{What it uses.} \todo{which properties of the latent space Palantir relies on}

\paragraph{Waddington-OT.}
Waddington-OT \cite{schiebinger2019optimal} links cells of consecutive days by optimal transport.
The cost of moving a cell to another is the squared Euclidean distance between them in the
latent space, so the cost matrix is the only part that depends on the latent space. The
transport is entropy-regularized (parameter $\varepsilon$) and unbalanced (parameter $\tau$),
because cells divide and die between days. Growth rates are estimated from two KEGG pathways
for mouse \cite{kanehisa2000kegg}: cell cycle (mmu04110) for proliferation and apoptosis
(mmu04210) for death. The same two pathways were used to estimate growth on these data in
PRESCIENT \cite{yeo2021prescient}. Of the 157 cell cycle genes, 154 are present in the data, and
all 136 apoptosis genes are present (four genes appear under older symbols: Casc5, Sgol1, Skp1a,
Sept4). Five genes (Atm, Gadd45a, Gadd45b, Gadd45g, Trp53) belong to both pathways. The gene
lists were retrieved from KEGG on 6 November 2026. We compute the growth rates once and use the same values and the same $\varepsilon$ and
$\tau$ for all latent spaces. The problem is solved with moscot \cite{klein2025moscot}.
\todo{values of $\varepsilon$ and $\tau$; how scores are turned into growth rates}

\emph{What it uses.} 
WOT assumes Markov nature of the data, requieres enough density in time, so it wouldn't need to do too big 
jumps in space, and it strongly depends on right hyperparameters.

\subsection{Latent spaces}
\label{sec:latent}
For each latent space we state a general assumption and, where a trajectory method is expected to
react differently, an assumption for that method.

\paragraph{PCA.}
Normalized counts were log-transformed ($\log(1+x)$), the 2000 most variable genes were
selected with the Scanpy default (Seurat) criterion \cite{wolf2018scanpy}, each gene was scaled
to unit variance with values clipped at 10, and the first 50 principal components were
computed by exact eigendecomposition of the covariance matrix. All 130{,}887 cells were used.

\emph{Assumption.} \todo{what this latent space changes in the geometry and why it should or should not help}

\emph{By method.} \todo{only where a method is expected to react differently}

\paragraph{Smoothed PCA.}
A $k$-nearest-neighbour graph was built on the PCA coordinates by exact Euclidean
search, with Scanpy connectivity weights $C$. The latent space was smoothed by $t$ steps of
diffusion, $X_t = P^t X$ with $P = D^{-1}(C+I)$, where $D$ is the diagonal matrix of row sums
of $C+I$. We use $t \in \{1, 3, 10\}$. This is a graph filter without learned weights and serves
as the control for graph autoencoders.

\emph{Assumption.} \todo{what this latent space changes in the geometry and why it should or should not help}

\emph{By method.} \todo{only where a method is expected to react differently}

\paragraph{MAGIC.}
MAGIC \cite{vandijk2018magic} imputes expression by diffusion on a cell--cell affinity graph.
\todo{parameters and how the latent is taken from the output}

\emph{Assumption.} \todo{what this latent space changes in the geometry and why it should or should not help}

\emph{By method.} \todo{only where a method is expected to react differently}

\subsection{Benchmark grid}
Table~\ref{tab:grid} lists the settings that are varied. Trained latent spaces are trained with 3 seeds. Since spread
between roots resulted to be small, we don't rerun model with different roots. A trained latent space without
an internal graph does not depend on $k$, so one model serves all values of $k$.

\begin{table}[ht]
\centering\small
\caption{Benchmark grid.}
\label{tab:grid}
\begin{tabular}{lp{0.62\linewidth}}
\toprule
Setting & Values \\
\midrule
Graph size $k$      & 10, 15, 30 \\
Latent space        & PCA; smoothed PCA ($t = 1, 3, 10$); MAGIC \todo{more} \\
Trajectory method   & Euclidean distance; DPT; Palantir; Slingshot; Waddington-OT \\
Training seed       & 0--2 (trained latent spaces) \\
\bottomrule
\end{tabular}
\end{table}

\section{Results}

\todo{results}

\bibliographystyle{plain}
\bibliography{references}

\appendix
\renewcommand{\thetable}{S\arabic{table}}
\renewcommand{\thefigure}{S\arabic{figure}}
\setcounter{table}{0}
\setcounter{figure}{0}

\section{Design checks}
All checks in this section use the development clones (1{,}760 clones, 26{,}106 sister pairs in
the two branches). The reported acc is the equal-weight mean over the neutrophil and monocyte
branches; intervals are 95\% bootstrap intervals over clones.

\subsection{Sensitivity to the root}
Five roots were drawn at random from the 100 undifferentiated day-2 cells closest to the centroid
of this population in PCA space. Table~\ref{tab:roots} shows acc for each root on PCA with
$k=15$. The standard deviation over roots is 0.008--0.012, and for Palantir acc ranges from 0.870
to 0.902 depending on the root. This is smaller than the bootstrap interval of a single method
(width about 0.06), but latent spaces are compared on the same clones, so most of the clone
sampling noise cancels in their differences while the effect of the root does not. We therefore
report acc averaged over the five roots for Euclidean distance, DPT and Palantir. Slingshot uses
only the cluster that contains the root, so it is run once.

\begin{table}[ht]
\centering\small
\caption{Accuracy on PCA ($k=15$) for five roots.}
\label{tab:roots}
\setlength{\tabcolsep}{3.5pt}
\begin{tabular}{lcccccccc}
\toprule
Method & 68853 & 16220 & 57472 & 48236 & 60816 & Mean & SD & CI width \\
\midrule
Euclidean distance       & 0.823 & 0.826 & 0.830 & 0.838 & 0.815 & 0.826 & 0.008 & 0.066 \\
DPT ($n_{\mathrm{dc}}=10$) & 0.870 & 0.873 & 0.878 & 0.891 & 0.865 & 0.876 & 0.010 & 0.064 \\
Palantir (20 comp.)      & 0.881 & 0.891 & 0.895 & 0.902 & 0.870 & 0.888 & 0.012 & 0.061 \\
\bottomrule
\end{tabular}
\end{table}

\subsection{Settings of the trajectory methods}
Settings were compared on PCA with $k=15$ (Table~\ref{tab:settings}; mean over the five roots,
interval averaged over roots). Differences between settings are at most 0.013, well within the
bootstrap interval. DPT uses 10 diffusion components (the Scanpy default). Palantir uses 20
components; this value is only an upper bound for its eigengap-based selection.
\todo{Slingshot}

\begin{table}[ht]
\centering\small
\caption{Accuracy on PCA ($k=15$) for different settings of the trajectory methods.}
\label{tab:settings}
\begin{tabular}{llcc}
\toprule
Method & Setting & acc & 95\% CI \\
\midrule
Euclidean distance & --- & 0.826 & 0.795--0.861 \\
\midrule
DPT & 5 diffusion components  & 0.869 & 0.838--0.903 \\
    & 10 diffusion components & 0.876 & 0.845--0.910 \\
    & 15 diffusion components & 0.875 & 0.844--0.909 \\
    & 20 diffusion components & 0.873 & 0.843--0.908 \\
    & 30 diffusion components & 0.870 & 0.839--0.904 \\
\midrule
Palantir & 5 components  & 0.875 & 0.849--0.901 \\
         & 10 components & 0.882 & 0.853--0.916 \\
         & 20 components & 0.888 & 0.860--0.920 \\
\bottomrule
\end{tabular}
\end{table}

\subsection{Separate k for frontend and backend graphs}
In the main grid the graph inside a frontend and the graph of the trajectory method use the same
$k$. To check this choice, smoothed PCA ($t=3$) was built with $k_f \in \{10, 15, 30\}$ and DPT
was run with $k_b \in \{10, 15, 30\}$ (root 68853; Table~\ref{tab:kcross}). Equal $k$ is never
worse than unequal $k$ beyond the bootstrap interval. The only large difference is $k_f=30$,
$k_b=10$, where acc drops to 0.701 (0.556 in the neutrophil branch).

\begin{table}[ht]
\centering\small
\caption{Accuracy of DPT on smoothed PCA ($t=3$) for frontend graph size $k_f$ (rows) and
backend graph size $k_b$ (columns). Intervals are about $\pm 0.03$ except for $k_f=30$, $k_b=10$
(0.659--0.756).}
\label{tab:kcross}
\begin{tabular}{lccc}
\toprule
$k_f$ \textbackslash{} $k_b$ & 10 & 15 & 30 \\
\midrule
10 & \textbf{0.862} & 0.864 & 0.863 \\
15 & 0.855 & \textbf{0.866} & 0.865 \\
30 & 0.701 & 0.859 & \textbf{0.866} \\
\bottomrule
\end{tabular}
\end{table}

\end{document}
